\subsection{AI Systems Recognizing and Responding to Moral Emotions}
\label{sec:5.2.3}
Some emotional states carry moral information: the guilt in a person's voice when admitting a mistake, the distress that signals a boundary has been crossed. Reading those uses the ordinary input channels — facial expression, vocal pattern, physiological signal — pointed at a narrower target. A healthcare AI system that detects anxiety in a patient's tone and adjusts its response accordingly is doing the recognition job section~\ref{sec:2.3.2} describes; a system that detects the same anxiety and treats it as a reason to weight the patient's stated preference more heavily in a treatment decision is doing the job this section is about.

The second half of the job is producing something in response — not necessarily a felt state, but a modeled one accurate enough to inform a decision. An autonomous vehicle that models the likely distress its emergency braking causes a passenger, and factors that into how sharply it brakes when a softer stop remains safe, is using a moral-emotion model without needing to have any emotion itself. The example is worth keeping and worth labelling: what it demonstrates is moral competence in section~\ref{sec:2.3}'s sense, recognizing a harm and letting the recognition bear on a decision. It demonstrates nothing about moral agency, and the vehicle would brake exactly as gently if the weighting were removed and reinstated by whoever maintains it. Most of what this chapter's methods deliver is competence of that kind, which is the right target for a vehicle and the wrong one for a bearer. Section~\ref{sec:5.3.1} takes up the deeper architectural question of how a fast, emotionally colored response and a slower deliberative one should combine in moral judgment generally; what belongs here is narrower — getting the emotional read right and using it honestly, prior to whatever combination rule section~\ref{sec:5.3} settles on.

Getting it right is harder than getting it fast. A system built to weigh moral emotion has to catch its own biases — the same distortion by race or gender that a human caregiver is vulnerable to does not disappear just because the caregiver is artificial — and it has to distinguish a genuine moral emotion from a performed one, since a display of distress calibrated to move an evaluator is exactly what a system trained to respond to distress will reward, whether or not the distress is real. Neither problem has a technical fix independent of the training data and feedback loop that shape it: ongoing exposure to varied, honestly labeled cases, and oversight positioned to catch a caretaker robot second-guessing a patient because of what the patient looks like and not what they need.
